\subsection{The Derivation IS the Computation}

The central insight of this framework is that \textbf{deriving a cell from partition operations IS computing cellular function}. There is no separation between:
\begin{itemize}
\item The mathematical derivation (partition operations on bounded phase space)
\item The physical system (the living cell)
\item The computation (cellular information processing)
\end{itemize}

These are three descriptions of the same underlying structure, just as oscillation $\equiv$ category $\equiv$ partition are three descriptions yielding identical entropy $S = k_B M \ln n$.

\begin{theorem}[Observation-Computation-Process Identity]
\label{thm:identity}
For any cellular process described by categorical morphism chain $\Phi_1 \circ \Phi_2 \circ \cdots \circ \Phi_k$:
\begin{equation}
\text{Observation} = \text{Computation} = \text{Process}
\end{equation}
The act of observing (projecting partition signatures), computing (completing trajectories), and physical process (cellular dynamics) are mathematically identical operations.
\end{theorem}

This identity resolves the apparent paradox of ``implementing'' cellular computation: we do not implement a model OF the cell---we run the same partition operations that constitute the cell. The Moon is derived from gravitational phase-lock networks; running those operations IS lunar dynamics. The cell is derived from categorical morphism chains; running those operations IS cellular function.

\subsection{The Poincar\'{e} Computing Paradigm}

Traditional computational biology simulates forward from initial conditions:
\begin{equation}
\mathbf{x}(t + \Delta t) = \mathbf{x}(t) + f(\mathbf{x}, t)\Delta t
\end{equation}

This requires $O(N \cdot T)$ operations for $N$ particles over $T$ time steps. For chaotic systems with Lyapunov exponent $\lambda$, precision requirements grow as $O(e^{\lambda T})$---computationally intractable for cellular timescales.

Poincar\'{e} computing inverts this paradigm: given boundary conditions and categorical constraints, valid trajectories are \emph{determined}, not computed step-by-step. The trajectory exists as a whole; we discover it through constraint propagation.

\begin{definition}[Poincar\'{e} Computing]
\label{def:poincare_computing}
Computational paradigm where trajectories are completed backward from observations through categorical constraints, rather than simulated forward from initial conditions using dynamical laws.
\end{definition}

\begin{theorem}[Complexity Reduction]
\label{thm:poincare_complexity}
Trajectory completion with $m$ constraints on $k$-trit precision requires $O(k \cdot m)$ operations, compared to $O(e^{\lambda T})$ for forward simulation of chaotic systems.
\end{theorem}

\begin{proof}
Each of $k$ trit positions requires checking $m$ constraints. Each constraint check is $O(1)$ for local constraints. Total: $O(k \cdot m)$. For cellular systems with $\lambda \sim 1$ s$^{-1}$ and characteristic times $T \sim 1$ ms, forward simulation requires $O(e^{1000})$ operations---effectively infinite. Trajectory completion requires $O(20 \times 3) = O(60)$ operations for typical enzymatic processes.
\end{proof}

\subsection{Position-Trajectory Duality}

The foundational insight enabling Poincar\'{e} computing is that in ternary S-entropy representation, position and trajectory are identical:

\begin{theorem}[Position-Trajectory Duality]
\label{thm:position_trajectory_duality}
For any $k$-trit string $\mathbf{t} = t_1 t_2 \cdots t_k$ where $t_i \in \{0, 1, 2\}$:
\begin{enumerate}
\item $\mathbf{t}$ addresses exactly one cell $C_{\mathbf{t}}$ in the $3^k$ partition of $\Sspace$
\item $\mathbf{t}$ encodes the unique trajectory from the origin to $C_{\mathbf{t}}$
\item The trajectory is the sequence of refinements: refine along axis $t_1$, then $t_2$, etc.
\end{enumerate}
The address IS the path.
\end{theorem}

This duality eliminates the von Neumann separation between data (position) and instructions (trajectory). In cellular computing, they are identical.

\subsection{Primitive Operations}

Three primitive operations replace Boolean AND, OR, NOT as computational foundations:

\begin{definition}[\textsc{Project}]
\label{def:project}
Extract categorical coordinate from physical observable:
\begin{equation}
\textsc{Project}: \Omega_{\text{physical}} \to \Sspace
\end{equation}
Given physical state $\Psi$, returns S-entropy coordinates $(\Sk, \St, \Se)$.

Implementation:
\begin{equation}
\textsc{Project}_i(\mathbf{t}) = t_i
\end{equation}
Extracts the $i$-th trit, revealing refinement along one axis.
\end{definition}

\begin{definition}[\textsc{Complete}]
\label{def:complete}
Given partial trajectory and constraints, return valid completion:
\begin{equation}
\textsc{Complete}: (\mathbf{t}_{1:j}, \mathcal{C}) \to \mathbf{t}_{j+1:k}
\end{equation}
This is the core operation of Poincar\'{e} computing. Given boundary conditions $\mathbf{t}_0$ and $\mathbf{t}_f$, plus categorical constraints $\mathcal{C}$, returns the unique trajectory satisfying all constraints.

When multiple completions exist, returns the set of valid completions.
\end{definition}

\begin{definition}[\textsc{Compose}]
\label{def:compose}
Concatenate two trajectories:
\begin{equation}
\textsc{Compose}: (\mathbf{t}^{(1)}, \mathbf{t}^{(2)}) \to \mathbf{t}^{(1)} \cdot \mathbf{t}^{(2)}
\end{equation}
The endpoint of $\mathbf{t}^{(1)}$ must match the startpoint of $\mathbf{t}^{(2)}$ in categorical space.
\end{definition}

\begin{proposition}[Closure]
The operations \textsc{Project}, \textsc{Complete}, and \textsc{Compose} are closed over the space of ternary trajectories: any composition of these operations yields a valid ternary trajectory.
\end{proposition}

\subsection{Backward Completion Algorithm}

\begin{definition}[Backward Completion]
\label{def:backward_completion}
Given final state $\mathbf{t}_f$ and constraints $\mathcal{C}$, propagate constraints backward:
\begin{enumerate}
\item Initialize: $\mathbf{t}^{(k)} = \mathbf{t}_f$
\item For $i = k-1$ down to $1$:
\begin{enumerate}
\item Compute valid predecessors: $\mathcal{P}_i = \{t : \mathcal{C}(t, \mathbf{t}^{(i+1:k)}) = \text{valid}\}$
\item Select: $t_i \in \mathcal{P}_i$ (unique if constraints are tight)
\end{enumerate}
\item Return: $\mathbf{t}^* = t_1 \cdots t_k$
\end{enumerate}
\end{definition}

\begin{theorem}[Backward Completion Complexity]
Backward completion with $m$ constraints over $k$ trit positions requires exactly $k \cdot m$ constraint evaluations.
\end{theorem}

\subsection{Constraint Types}

Four constraint types govern trajectory validity:

\begin{enumerate}
\item \textbf{Boundary constraints:} Fix initial/final states
\begin{equation}
\mathcal{C}_{\text{boundary}}: \mathbf{t}_0 = (0,0,0,\ldots), \quad \mathbf{t}_f = (2,2,2,\ldots)
\end{equation}

\item \textbf{Aperture constraints:} Require passage through categorical aperture $\mathcal{A}$
\begin{equation}
\mathcal{C}_{\mathcal{A}}(\mathbf{t}) = \begin{cases}
\text{valid} & \exists i: C_{t_1 \cdots t_i} \cap \mathcal{A} \neq \emptyset \\
\text{invalid} & \text{otherwise}
\end{cases}
\end{equation}

\item \textbf{Continuity constraints:} Adjacent trits satisfy $|t_{i+1} - t_i| \leq 1$

\item \textbf{Conservation constraints:} Certain trit patterns sum to fixed values (charge neutrality, energy conservation)
\end{enumerate}

\subsection{Data Structures}

\subsubsection{Ternary Tree}

Hierarchical partition structure storing categorical states:

\begin{verbatim}
struct TernaryNode {
    depth: u32,
    coordinate: (f64, f64, f64),  // (S_k, S_t, S_e)
    children: [Option<Box<TernaryNode>>; 3],
    constraints: Vec<Constraint>,
    valid: bool,
}
\end{verbatim}

Tree operations:
\begin{itemize}
\item \texttt{insert(trit\_string)}: $O(\log_3 N)$
\item \texttt{lookup(trit\_string)}: $O(\log_3 N)$
\item \texttt{validate(constraints)}: $O(k \cdot m)$
\item \texttt{complete(boundary, constraints)}: $O(k \cdot m)$
\end{itemize}

\subsubsection{Oscillator Network}

Phase-locked oscillator array for harmonic coincidence:

\begin{verbatim}
struct OscillatorNetwork {
    oscillators: Vec<Oscillator>,
    coupling_matrix: Matrix<f64>,
    phase_lock_threshold: f64,
}

struct Oscillator {
    frequency: f64,           // Hz
    phase: f64,               // radians
    amplitude: f64,
    coupling_strength: f64,
}
\end{verbatim}

Network operations:
\begin{itemize}
\item \texttt{detect\_coincidence()}: Returns beat frequencies where phases align
\item \texttt{categorical\_access(S\_coord)}: Maps S-entropy coordinate to oscillator state
\item \texttt{phase\_lock(oscillators)}: Synchronize oscillator subset
\end{itemize}

\subsubsection{Constraint Graph}

Dependency graph for constraint propagation:

\begin{verbatim}
struct ConstraintGraph {
    nodes: Vec<ConstraintNode>,
    edges: Vec<(usize, usize)>,  // dependency edges
}

struct ConstraintNode {
    constraint_type: ConstraintType,
    trit_positions: Vec<usize>,
    evaluate: fn(&TritString) -> bool,
}
\end{verbatim}

\subsection{Cellular Partition Language (CPL)}

A declarative language for specifying cellular computations:

\subsubsection{Syntax}

\begin{verbatim}
# Define S-entropy space
space S = [0,1]^3

# Define categorical aperture (enzyme active site)
aperture CA_II {
    geometry: tetrahedral
    center: (0.5, 0.5, 0.5)
    width: 0.01
    pattern: "012"  # zinc passage trit sequence
}

# Define boundary conditions
initial = project(CO2, r=5Å)  # substrate
final = project(HCO3-, r=5Å)  # product

# Define constraints
constraints {
    charge_neutrality: sum(q_i) == 0
    energy_conservation: |E_f - E_i| < kT
    categorical_coherence: R > 0.7
}

# Complete trajectory
trajectory = complete(
    from: initial,
    to: final,
    through: [CA_II],
    satisfying: constraints,
    precision: 20 trits
)

# Extract observables
categorical_distance = |trajectory|
transitions = count(trajectory, pattern="012")
turnover = omega_0 * exp(-categorical_distance / k_B)
\end{verbatim}

\subsubsection{Semantics}

\begin{enumerate}
\item \texttt{space}: Declares S-entropy coordinate space
\item \texttt{aperture}: Defines geometric constraint (enzyme, channel, etc.)
\item \texttt{project}: Maps physical observables to S-entropy coordinates
\item \texttt{constraints}: Specifies physical validity constraints
\item \texttt{complete}: Core operation---backward propagation from boundary conditions
\item Observable extraction: Derive measurable quantities from completed trajectory
\end{enumerate}

\subsubsection{No Dynamical Laws}

CPL contains no force fields, rate equations, or differential equations. Trajectories emerge from constraint satisfaction in S-entropy space. Laws are \emph{derived} from constraints through geometry:
\begin{itemize}
\item Euler-Lagrange equations: Geodesic constraint
\item Fick's law: Diffusive aperture constraint
\item Michaelis-Menten: Enzymatic aperture constraint
\end{itemize}

\subsection{Interface Specification}

\subsubsection{Input Format}

\begin{verbatim}
{
    "query_type": "trajectory_completion",
    "initial_state": {
        "S_k": 0.1,
        "S_t": 0.1,
        "S_e": 0.1
    },
    "final_state": {
        "S_k": 0.9,
        "S_t": 0.9,
        "S_e": 0.9
    },
    "apertures": [
        {
            "name": "enzyme_active_site",
            "center": [0.5, 0.5, 0.5],
            "width": 0.01,
            "pattern": "012"
        }
    ],
    "constraints": [
        {"type": "charge_neutrality"},
        {"type": "energy_conservation", "tolerance": 1e-3},
        {"type": "categorical_coherence", "threshold": 0.7}
    ],
    "precision": 20
}
\end{verbatim}

\subsubsection{Output Format}

\begin{verbatim}
{
    "status": "completed",
    "trajectory": {
        "trit_string": "00000000001201222222",
        "length": 20,
        "categorical_distance": 1,
        "transitions": [
            {"position": 10, "pattern": "012", "type": "aperture_traversal"}
        ]
    },
    "observables": {
        "path_length": 20,
        "turnover_rate": 1e6,
        "constraint_satisfaction": {
            "charge_neutrality": true,
            "energy_conservation": true,
            "categorical_coherence": 0.85
        }
    },
    "metadata": {
        "computation_time_ns": 1250,
        "constraint_checks": 60,
        "oscillator_coincidences": 42
    }
}
\end{verbatim}

\subsection{Hardware Mapping}

\subsubsection{Consumer Hardware Components}

\begin{table}[h]
\centering
\begin{tabular}{lcc}
\toprule
Component & Frequency & Role \\
\midrule
CPU clock & $\sim 3$ GHz & Base oscillator \\
GPU clock & $\sim 2$ GHz & Parallel coincidence detection \\
Memory controller & $\sim 3$ GHz & Phase reference \\
USB controller & $\sim 480$ MHz & Low-frequency harmonics \\
Audio DAC & $\sim 192$ kHz & Fine timing \\
Network PHY & $\sim 125$ MHz & Independent reference \\
PCIe clock & $\sim 100$ MHz & Synchronization \\
RTC crystal & $32.768$ kHz & Long-term stability \\
\bottomrule
\end{tabular}
\caption{Consumer hardware oscillators for harmonic coincidence network}
\end{table}

\subsubsection{Coincidence Detection}

Beat frequencies from $N = 8$ oscillators with integer coefficients $n_i \in [-10^6, 10^6]$:
\begin{equation}
\omega_{\text{beat}} = \left|\sum_{i=1}^N n_i \omega_i\right|
\end{equation}

Coincidence occurs when multiple beat frequency paths yield identical values:
\begin{equation}
\omega_{\text{beat}}^{(1)} = \omega_{\text{beat}}^{(2)} = \cdots = \omega_{\text{beat}}^{(M)}
\end{equation}

Each coincidence maps to a categorical coordinate through:
\begin{equation}
(\Sk, \St, \Se) = \text{decode}(\omega_{\text{beat}})
\end{equation}

\subsection{Implementation Phases}

\subsubsection{Phase 1: Python Validation}

Rapid prototyping to validate:
\begin{itemize}
\item Ternary tree operations
\item Backward completion algorithm
\item Constraint propagation
\item CPL parser and interpreter
\item CA II trajectory reproduction
\end{itemize}

Target: Reproduce $10^9\times$ speedup over molecular dynamics for enzymatic trajectories.

\subsubsection{Phase 2: Rust Implementation}

Production implementation with:
\begin{itemize}
\item Zero-cost abstractions for ternary operations
\item Lock-free concurrent constraint evaluation
\item Hardware oscillator interface (via OS timing APIs)
\item SIMD-optimized coincidence detection
\item Memory-mapped partition trees
\end{itemize}

Target: Real-time cellular trajectory completion at $10^6$ queries/second.

\subsubsection{Phase 3: Hardware Integration}

Direct hardware oscillator access:
\begin{itemize}
\item FPGA-based coincidence detection
\item Custom oscillator arrays
\item Phase-locked loop integration
\item Trans-Planckian timing validation
\end{itemize}

Target: Categorical access at $\tau_{\text{meas}} = 10^{-66}$ s effective resolution.
